% Options for packages loaded elsewhere
\PassOptionsToPackage{unicode}{hyperref}
\PassOptionsToPackage{hyphens}{url}
\documentclass[
]{article}
\usepackage{xcolor}
\usepackage[margin=1in]{geometry}
\usepackage{amsmath,amssymb}
\setcounter{secnumdepth}{-\maxdimen} % remove section numbering
\usepackage{iftex}
\ifPDFTeX
  \usepackage[T1]{fontenc}
  \usepackage[utf8]{inputenc}
  \usepackage{textcomp} % provide euro and other symbols
\else % if luatex or xetex
  \usepackage{unicode-math} % this also loads fontspec
  \defaultfontfeatures{Scale=MatchLowercase}
  \defaultfontfeatures[\rmfamily]{Ligatures=TeX,Scale=1}
\fi
\usepackage{lmodern}
\ifPDFTeX\else
  % xetex/luatex font selection
\fi
% Use upquote if available, for straight quotes in verbatim environments
\IfFileExists{upquote.sty}{\usepackage{upquote}}{}
\IfFileExists{microtype.sty}{% use microtype if available
  \usepackage[]{microtype}
  \UseMicrotypeSet[protrusion]{basicmath} % disable protrusion for tt fonts
}{}
\makeatletter
\@ifundefined{KOMAClassName}{% if non-KOMA class
  \IfFileExists{parskip.sty}{%
    \usepackage{parskip}
  }{% else
    \setlength{\parindent}{0pt}
    \setlength{\parskip}{6pt plus 2pt minus 1pt}}
}{% if KOMA class
  \KOMAoptions{parskip=half}}
\makeatother
\usepackage{longtable,booktabs,array}
\newcounter{none} % for unnumbered tables
\usepackage{calc} % for calculating minipage widths
% Correct order of tables after \paragraph or \subparagraph
\usepackage{etoolbox}
\makeatletter
\patchcmd\longtable{\par}{\if@noskipsec\mbox{}\fi\par}{}{}
\makeatother
% Allow footnotes in longtable head/foot
\IfFileExists{footnotehyper.sty}{\usepackage{footnotehyper}}{\usepackage{footnote}}
\makesavenoteenv{longtable}
\usepackage{graphicx}
\makeatletter
\newsavebox\pandoc@box
\newcommand*\pandocbounded[1]{% scales image to fit in text height/width
  \sbox\pandoc@box{#1}%
  \Gscale@div\@tempa{\textheight}{\dimexpr\ht\pandoc@box+\dp\pandoc@box\relax}%
  \Gscale@div\@tempb{\linewidth}{\wd\pandoc@box}%
  \ifdim\@tempb\p@<\@tempa\p@\let\@tempa\@tempb\fi% select the smaller of both
  \ifdim\@tempa\p@<\p@\scalebox{\@tempa}{\usebox\pandoc@box}%
  \else\usebox{\pandoc@box}%
  \fi%
}
% Set default figure placement to htbp
\def\fps@figure{htbp}
\makeatother
\setlength{\emergencystretch}{3em} % prevent overfull lines
\providecommand{\tightlist}{%
  \setlength{\itemsep}{0pt}\setlength{\parskip}{0pt}}
\usepackage{bookmark}
\IfFileExists{xurl.sty}{\usepackage{xurl}}{} % add URL line breaks if available
\urlstyle{same}
\hypersetup{
  pdftitle={In-Hospital Cardiac Arrest \& ICU Outcomes Audit},
  pdfauthor={Hussein Kassir},
  hidelinks,
  pdfcreator={LaTeX via pandoc}}

\title{In-Hospital Cardiac Arrest \& ICU Outcomes Audit}
\usepackage{etoolbox}
\makeatletter
\providecommand{\subtitle}[1]{% add subtitle to \maketitle
  \apptocmd{\@title}{\par {\large #1 \par}}{}{}
}
\makeatother
\subtitle{Modelled on ICNARC National Cardiac Arrest Audit --- MIMIC-IV
Demo Dataset}
\author{Hussein Kassir}
\date{2026-06-09}

\begin{document}
\maketitle

{
\setcounter{tocdepth}{2}
\tableofcontents
}
\section{Executive Summary}\label{executive-summary}

This report presents findings from an audit of in-hospital cardiac
arrest and ICU outcomes, modelled on the methodology of ICNARC's
National Cardiac Arrest Audit (NCAA) and Case Mix Programme (CMP). The
audit was conducted using the MIMIC-IV Clinical Database Demo (v2.2), an
openly available dataset of 100 de-identified critical care patients
from Beth Israel Deaconess Medical Center, Boston.

\textbf{Key findings:}

\begin{itemize}
\tightlist
\item
  One cardiac arrest case was identified in the demo dataset (ICD-10:
  I462), admitted as an emergency to the Coronary Care Unit, with
  in-hospital death confirmed.
\item
  Across the full ICU cohort of 100 patients (140 stays), overall
  in-hospital mortality was \textbf{13\%}.
\item
  The Coronary Care Unit recorded the highest mortality rate at
  \textbf{38.5\%}, followed by the Neuro Surgical ICU at
  \textbf{33.3\%}.
\item
  ICU length of stay was the only statistically significant predictor of
  in-hospital death in the risk model (OR 1.30, 95\% CI 1.12--1.58,
  p=0.002).
\item
  51 data quality discrepancies were identified and logged, of which 9
  were rated High severity and 42 Medium severity.
\end{itemize}

\begin{quote}
\textbf{Note:} This audit was developed as a portfolio project to
demonstrate clinical audit data management and analytical skills. The
MIMIC-IV Demo dataset contains 100 patients and results are not
statistically powered for clinical inference. The pipeline is designed
to scale to the full MIMIC-IV dataset.
\end{quote}

\section{Background \& Methods}\label{background-methods}

\subsection{Audit Scope}\label{audit-scope}

This audit was designed to mirror the structure and objectives of
ICNARC's National Cardiac Arrest Audit (NCAA), which monitors the
incidence, care, and outcomes of in-hospital cardiac arrest across NHS
hospitals in England, Wales, and Northern Ireland. A secondary analysis
of broader ICU outcomes was conducted across the full cohort, modelled
on ICNARC's Case Mix Programme (CMP).

\subsection{Dataset}\label{dataset}

The \textbf{MIMIC-IV Clinical Database Demo (v2.2)} was used as the data
source. MIMIC-IV is a large, freely available database of de-identified
health data from patients admitted to critical care units at Beth Israel
Deaconess Medical Center between 2008 and 2019. The demo version
contains 100 patients and is openly available via PhysioNet without
credentialing requirements.

\subsection{Cohort Definition}\label{cohort-definition}

\textbf{Cardiac arrest cohort:}

Patients were included if they had a recorded diagnosis of cardiac
arrest using the following codes:

\begin{itemize}
\tightlist
\item
  ICD-10: I46, I460, I461, I462, I469
\item
  ICD-9: 4275
\end{itemize}

Patients were excluded if aged under 18 or flagged with a high-severity
data quality discrepancy.

\textbf{Full ICU cohort:}

All patients with at least one ICU stay in the dataset were included in
the secondary analysis, subject to the same age and QC exclusion
criteria.

\subsection{Outcome Measures}\label{outcome-measures}

The primary outcome was \textbf{in-hospital mortality}, defined as
\texttt{hospital\_expire\_flag\ =\ 1} in the admissions table. Secondary
outcomes included ICU and hospital length of stay, care unit
distribution, and admission type breakdown.

\subsection{Analytical Pipeline}\label{analytical-pipeline}

All data processing, quality control, cohort definition, and analysis
were conducted in R (v4.x) using a reproducible, version-controlled
pipeline maintained on GitHub. Data quality checks were modelled on
ICNARC's discrepancy resolution process, with all flagged records logged
with a severity rating and resolution status.

\section{Data Quality Summary}\label{data-quality-summary}

\subsection{Overview}\label{overview}

Data quality checks were conducted across all four core tables:
\texttt{patients}, \texttt{admissions}, \texttt{diagnoses\_icd}, and
\texttt{icustays}. Three categories of checks were applied:
completeness, plausibility, and consistency. All flagged records were
logged in a structured discrepancy log with a severity rating and
resolution status, replicating the query-raising process used in live
clinical audits.

\subsection{Completeness}\label{completeness}

{\def\LTcaptype{none} % do not increment counter
\begin{longtable}[]{@{}
  >{\raggedright\arraybackslash}p{(\linewidth - 6\tabcolsep) * \real{0.1048}}
  >{\raggedright\arraybackslash}p{(\linewidth - 6\tabcolsep) * \real{0.2095}}
  >{\raggedright\arraybackslash}p{(\linewidth - 6\tabcolsep) * \real{0.0952}}
  >{\raggedright\arraybackslash}p{(\linewidth - 6\tabcolsep) * \real{0.5905}}@{}}
\toprule\noalign{}
\begin{minipage}[b]{\linewidth}\raggedright
Table
\end{minipage} & \begin{minipage}[b]{\linewidth}\raggedright
Variable
\end{minipage} & \begin{minipage}[b]{\linewidth}\raggedright
\% Missing
\end{minipage} & \begin{minipage}[b]{\linewidth}\raggedright
Interpretation
\end{minipage} \\
\midrule\noalign{}
\endhead
\bottomrule\noalign{}
\endlastfoot
patients & dod & 69\% & Expected --- only populated for deaths \\
admissions & deathtime & 94.5\% & Expected --- only populated for
in-hospital deaths \\
admissions & edregtime / edouttime & 33.8\% & Expected --- not all
patients admitted via emergency department \\
admissions & discharge\_location & 15.3\% & Flagged --- unknown
discharge destination, audit-relevant \\
admissions & marital\_status & 4.4\% & Minor gap, low audit relevance \\
\end{longtable}
}

\subsection{Plausibility}\label{plausibility}

No implausible values were identified for age, ICU length of stay, or
admission/discharge dates. All 100 patients were adults (≥18), all ICU
lengths of stay were positive, and no discharge dates preceded admission
dates.

\subsection{Consistency}\label{consistency}

Four records showed death time recorded after discharge time by a few
hours. Investigation revealed this reflects a known MIMIC data
convention where discharge time is recorded as midnight (00:00:00) and
death time is logged precisely later that evening. These were documented
as known artefacts and excluded from the discrepancy log.

Nine records were flagged for ICU stay timestamps falling outside the
hospital admission window by more than six hours, likely reflecting
date-shift artefacts in the MIMIC de-identification process.

\subsection{Discrepancy Log Summary}\label{discrepancy-log-summary}

{\def\LTcaptype{none} % do not increment counter
\begin{longtable}[]{@{}llr@{}}
\toprule\noalign{}
Check Type & Severity & Records Flagged \\
\midrule\noalign{}
\endhead
\bottomrule\noalign{}
\endlastfoot
Consistency & High & 9 \\
Completeness & Medium & 42 \\
\end{longtable}
}

In a live audit setting, High severity discrepancies would be queried
back to the submitting unit for clarification or correction before
inclusion in analysis. Medium severity records would be flagged for
review and included with a data quality caveat.

\section{Cardiac Arrest Cohort}\label{cardiac-arrest-cohort}

\subsection{Cohort identification}\label{cohort-identification}

{\def\LTcaptype{none} % do not increment counter
\begin{longtable}[]{@{}lr@{}}
\toprule\noalign{}
Step & N \\
\midrule\noalign{}
\endhead
\bottomrule\noalign{}
\endlastfoot
All MIMIC-IV Demo patients & 100 \\
Cardiac arrest ICD code identified & 1 \\
Excluded --- under 18 & 0 \\
Excluded --- high severity QC flag & 0 \\
Final analysis cohort & 1 \\
\end{longtable}
}

\subsection{Patient Profile}\label{patient-profile}

{\def\LTcaptype{none} % do not increment counter
\begin{longtable}[]{@{}ll@{}}
\toprule\noalign{}
Variable & Value \\
\midrule\noalign{}
\endhead
\bottomrule\noalign{}
\endlastfoot
Gender & FALSE \\
Age & 89 \\
Admission Type & EW EMER. \\
Admission Location & EMERGENCY ROOM \\
Care Unit & Coronary Care Unit (CCU) \\
ICU LOS (days) & 4.90715277777778 \\
Hospital LOS (days) & 4.82916666666667 \\
In-Hospital Death & Yes \\
ICD Code & I462 \\
\end{longtable}
}

\subsection{Findings \& Limitations}\label{findings-limitations}

The single cardiac arrest case identified was a female patient aged 89,
admitted as an emergency via the emergency department and managed in the
Coronary Care Unit. In-hospital death was confirmed, with an ICU length
of stay of 4.9 days and hospital length of stay of 4.8 days. The arrest
was coded as ICD-10 I462 (cardiac arrest due to underlying cardiac
condition).

\begin{quote}
\textbf{Limitation:} The MIMIC-IV Demo dataset contains only 100
patients, yielding a single cardiac arrest case. Outcome statistics for
this cohort are not meaningful at this scale. The identification,
validation, and reporting pipeline is fully operational and designed to
scale to the full MIMIC-IV dataset, which contains over 300,000 hospital
admissions.
\end{quote}

\section{Full ICU Cohort Analysis}\label{full-icu-cohort-analysis}

\subsection{Overview}\label{overview-1}

A secondary analysis was conducted across all 100 patients with at least
one ICU stay in the dataset, yielding 140 ICU stays. This analysis
mirrors the Case Mix Programme's approach to benchmarking critical care
outcomes across participating units.

\subsection{In-Hospital Mortality}\label{in-hospital-mortality}

{\def\LTcaptype{none} % do not increment counter
\begin{longtable}[]{@{}ll@{}}
\toprule\noalign{}
Metric & Value \\
\midrule\noalign{}
\endhead
\bottomrule\noalign{}
\endlastfoot
Total patients & 100 \\
In-hospital deaths & 13 \\
Survivors & 87 \\
Mortality rate & 13\% \\
\end{longtable}
}

\subsection{ICU Length of Stay}\label{icu-length-of-stay}

{\def\LTcaptype{none} % do not increment counter
\begin{longtable}[]{@{}ll@{}}
\toprule\noalign{}
Metric & Value \\
\midrule\noalign{}
\endhead
\bottomrule\noalign{}
\endlastfoot
Mean ICU LOS (days) & 3.68 \\
Median ICU LOS (days) & 2.16 \\
Minimum ICU LOS (days) & 0.02 \\
Maximum ICU LOS (days) & 20.5 \\
\end{longtable}
}

The median ICU length of stay was 2.16 days, with considerable variation
(range 0.02--20.5 days), reflecting the heterogeneity of case mix
typical of a mixed critical care population.

\subsection{Breakdown by Care Unit}\label{breakdown-by-care-unit}

{\def\LTcaptype{none} % do not increment counter
\begin{longtable}[]{@{}
  >{\raggedright\arraybackslash}p{(\linewidth - 8\tabcolsep) * \real{0.5104}}
  >{\raggedleft\arraybackslash}p{(\linewidth - 8\tabcolsep) * \real{0.1042}}
  >{\raggedleft\arraybackslash}p{(\linewidth - 8\tabcolsep) * \real{0.0729}}
  >{\raggedleft\arraybackslash}p{(\linewidth - 8\tabcolsep) * \real{0.1458}}
  >{\raggedleft\arraybackslash}p{(\linewidth - 8\tabcolsep) * \real{0.1667}}@{}}
\toprule\noalign{}
\begin{minipage}[b]{\linewidth}\raggedright
Care Unit
\end{minipage} & \begin{minipage}[b]{\linewidth}\raggedleft
ICU Stays
\end{minipage} & \begin{minipage}[b]{\linewidth}\raggedleft
Deaths
\end{minipage} & \begin{minipage}[b]{\linewidth}\raggedleft
Mortality (\%)
\end{minipage} & \begin{minipage}[b]{\linewidth}\raggedleft
Mean LOS (days)
\end{minipage} \\
\midrule\noalign{}
\endhead
\bottomrule\noalign{}
\endlastfoot
Medical Intensive Care Unit (MICU) & 29 & 6 & 20.7 & 4.02 \\
Surgical Intensive Care Unit (SICU) & 29 & 3 & 10.3 & 2.83 \\
Cardiac Vascular Intensive Care Unit (CVICU) & 25 & 0 & 0.0 & 2.22 \\
Medical/Surgical Intensive Care Unit (MICU/SICU) & 23 & 3 & 13.0 &
4.25 \\
Trauma SICU (TSICU) & 16 & 2 & 12.5 & 4.07 \\
Coronary Care Unit (CCU) & 13 & 5 & 38.5 & 5.43 \\
Neuro Surgical Intensive Care Unit (Neuro SICU) & 3 & 1 & 33.3 & 6.33 \\
Neuro Intermediate & 1 & 0 & 0.0 & 0.87 \\
Neuro Stepdown & 1 & 0 & 0.0 & 7.70 \\
\end{longtable}
}

The Coronary Care Unit recorded the highest mortality rate at 38.5\%,
followed by the Neuro Surgical ICU at 33.3\%. The Cardiac Vascular ICU
recorded zero deaths across 25 stays, likely reflecting a higher
proportion of elective surgical admissions with lower baseline risk.

\subsection{Breakdown by Age Band}\label{breakdown-by-age-band}

{\def\LTcaptype{none} % do not increment counter
\begin{longtable}[]{@{}lrrr@{}}
\toprule\noalign{}
Age Band & Patients & Deaths & Mortality (\%) \\
\midrule\noalign{}
\endhead
\bottomrule\noalign{}
\endlastfoot
40-59 & 33 & 2 & 6.1 \\
60-79 & 42 & 9 & 21.4 \\
80 and over & 16 & 2 & 12.5 \\
Under 40 & 9 & 0 & 0.0 \\
\end{longtable}
}

Mortality was highest in the 60--79 age band at 21.4\%. The apparent
lower mortality in the 80 and over group (12.5\%) may reflect survivor
bias --- only the fittest older patients are admitted to ICU --- and
should be interpreted with caution given the small sample size.

\subsection{Breakdown by Admission
Type}\label{breakdown-by-admission-type}

{\def\LTcaptype{none} % do not increment counter
\begin{longtable}[]{@{}lrrr@{}}
\toprule\noalign{}
Admission Type & Patients & Deaths & Mortality (\%) \\
\midrule\noalign{}
\endhead
\bottomrule\noalign{}
\endlastfoot
EW EMER. & 46 & 6 & 13.0 \\
URGENT & 22 & 4 & 18.2 \\
SURGICAL SAME DAY ADMISSION & 14 & 0 & 0.0 \\
OBSERVATION ADMIT & 11 & 2 & 18.2 \\
ELECTIVE & 5 & 0 & 0.0 \\
DIRECT EMER. & 2 & 1 & 50.0 \\
\end{longtable}
}

Elective and surgical same-day admissions recorded zero mortality,
consistent with their lower acuity. Direct emergency admissions had the
highest mortality at 50\%, though this is based on only two patients and
should not be interpreted as a reliable estimate.

\section{Risk-Adjusted Outcomes}\label{risk-adjusted-outcomes}

\subsection{Logistic Regression Model}\label{logistic-regression-model}

A logistic regression model was fitted to identify predictors of
in-hospital mortality across the full ICU cohort. Predictor variables
were selected to reflect those used in ICNARC's risk prediction models:
age, sex, admission type (emergency vs non-emergency), and ICU length of
stay.

{\def\LTcaptype{none} % do not increment counter
\begin{longtable}[]{@{}
  >{\raggedright\arraybackslash}p{(\linewidth - 14\tabcolsep) * \real{0.2736}}
  >{\raggedleft\arraybackslash}p{(\linewidth - 14\tabcolsep) * \real{0.1038}}
  >{\raggedleft\arraybackslash}p{(\linewidth - 14\tabcolsep) * \real{0.1038}}
  >{\raggedleft\arraybackslash}p{(\linewidth - 14\tabcolsep) * \real{0.1132}}
  >{\raggedleft\arraybackslash}p{(\linewidth - 14\tabcolsep) * \real{0.0755}}
  >{\raggedleft\arraybackslash}p{(\linewidth - 14\tabcolsep) * \real{0.1038}}
  >{\raggedleft\arraybackslash}p{(\linewidth - 14\tabcolsep) * \real{0.1132}}
  >{\raggedright\arraybackslash}p{(\linewidth - 14\tabcolsep) * \real{0.1132}}@{}}
\toprule\noalign{}
\begin{minipage}[b]{\linewidth}\raggedright
Predictor
\end{minipage} & \begin{minipage}[b]{\linewidth}\raggedleft
Odds Ratio
\end{minipage} & \begin{minipage}[b]{\linewidth}\raggedleft
Std. Error
\end{minipage} & \begin{minipage}[b]{\linewidth}\raggedleft
Z Statistic
\end{minipage} & \begin{minipage}[b]{\linewidth}\raggedleft
P Value
\end{minipage} & \begin{minipage}[b]{\linewidth}\raggedleft
95\% CI Low
\end{minipage} & \begin{minipage}[b]{\linewidth}\raggedleft
95\% CI High
\end{minipage} & \begin{minipage}[b]{\linewidth}\raggedright
Significant
\end{minipage} \\
\midrule\noalign{}
\endhead
\bottomrule\noalign{}
\endlastfoot
Age (per year) & 1.038 & 0.023 & 1.600 & 0.110 & 0.994 & 1.091 & No \\
Emergency admission & 2.910 & 0.882 & 1.211 & 0.226 & 0.613 & 22.477 &
No \\
Female sex & 0.363 & 0.787 & -1.287 & 0.198 & 0.064 & 1.537 & No \\
ICU length of stay (per day) & 1.305 & 0.085 & 3.126 & 0.002 & 1.120 &
1.576 & Yes \\
\end{longtable}
}

\subsection{Key Findings}\label{key-findings}

\begin{itemize}
\tightlist
\item
  \textbf{ICU length of stay} was the only statistically significant
  predictor of in-hospital death (OR 1.30, 95\% CI 1.12--1.58, p=0.002).
  Each additional day in ICU was associated with a 30\% increase in the
  odds of in-hospital death, likely reflecting greater illness severity
  in patients with prolonged stays.
\item
  \textbf{Age} showed a positive but non-significant association with
  mortality (OR 1.04, p=0.11), consistent with the expected direction
  but underpowered at this sample size.
\item
  \textbf{Emergency admission} was associated with higher odds of death
  (OR 2.91) but with very wide confidence intervals (0.61--22.5),
  reflecting insufficient statistical power.
\item
  \textbf{Female sex} showed a non-significant protective association
  (OR 0.36, p=0.198).
\end{itemize}

\subsection{Observed vs Expected
Mortality}\label{observed-vs-expected-mortality}

{\def\LTcaptype{none} % do not increment counter
\begin{longtable}[]{@{}ll@{}}
\toprule\noalign{}
Metric & Value \\
\midrule\noalign{}
\endhead
\bottomrule\noalign{}
\endlastfoot
Observed deaths & 13 \\
Expected deaths & 13.0 \\
O/E ratio & 1.00 \\
\end{longtable}
}

The observed-to-expected (O/E) ratio of 1.00 indicates perfect
calibration of the model on this dataset. In a live audit context, an
O/E ratio substantially above 1.0 would indicate higher-than-expected
mortality and trigger a quality improvement review, while a ratio below
1.0 would suggest better-than-expected performance. This mirrors the
Risk-Adjusted Mortality Indicator (RAMI) used in ICNARC's Case Mix
Programme.

\begin{quote}
\textbf{Note:} The O/E ratio here reflects in-sample model fit rather
than external validation. In a scaled audit, the model would be trained
on a derivation dataset and validated on a separate cohort.
\end{quote}

\section{Limitations \&
Recommendations}\label{limitations-recommendations}

\subsection{Limitations}\label{limitations}

\textbf{Sample size:} The MIMIC-IV Demo dataset contains 100 patients,
yielding a single cardiac arrest case. All statistics derived from this
cohort should be interpreted as illustrative of the pipeline methodology
rather than clinically meaningful findings. The full MIMIC-IV dataset
contains over 300,000 hospital admissions and would provide adequate
statistical power for a fully powered audit.

\textbf{Single-centre data:} MIMIC-IV draws exclusively from Beth Israel
Deaconess Medical Center in Boston. A national audit such as ICNARC's
NCAA draws from hundreds of participating units across England, Wales,
and Northern Ireland, enabling meaningful inter-unit benchmarking.
Single-centre data cannot replicate this.

\textbf{ICD coding limitations:} Cardiac arrest identification relied
solely on ICD diagnostic codes. In clinical practice, coding
completeness varies across units and cardiac arrest may be undercoded,
particularly when it occurs as a complication rather than a primary
diagnosis. A live audit would supplement ICD coding with direct data
entry by clinical staff.

\textbf{Risk model variables:} The logistic regression model included
only four predictor variables. ICNARC's own risk prediction models
incorporate a much richer set of physiological variables collected at
ICU admission, including acute physiology scores, ventilation status,
and primary diagnosis category. These variables were not included here
due to the complexity of extracting and cleaning chartevents data within
the scope of this project.

\textbf{Date shifting:} MIMIC-IV applies date shifting to all timestamps
as part of its de-identification process. This introduced apparent
inconsistencies in some timestamp comparisons, particularly between ICU
stay times and hospital admission windows. These were identified through
QC checks and documented in the discrepancy log.

\subsection{Recommendations}\label{recommendations}

\begin{enumerate}
\def\labelenumi{\arabic{enumi}.}
\item
  \textbf{Scale to full MIMIC-IV dataset} --- applying this pipeline to
  the full dataset would yield a statistically powered cardiac arrest
  cohort and enable meaningful benchmarking analysis.
\item
  \textbf{Enrich the risk model} --- incorporating physiological
  variables from \texttt{chartevents} (e.g.~GCS, heart rate, blood
  pressure at ICU admission) would produce a more clinically credible
  risk adjustment model, closer to ICNARC's RAMI methodology.
\item
  \textbf{Automate discrepancy reporting} --- in a live audit setting,
  the discrepancy log could be extended to generate automated query
  letters to participating units, flagging specific records and
  requesting clarification or correction.
\item
  \textbf{Extend to NCAA variables} --- where derivable from MIMIC-IV,
  additional NCAA variables such as arrest rhythm (shockable vs
  non-shockable) and ROSC status could be extracted from clinical notes
  or procedure codes to more closely replicate the full NCAA dataset
  specification.
\end{enumerate}

\end{document}
